\chapter{空気を通した会話}
% full version: chapters/05_serocomputer.tex

\pencilfig{5}{\pencilart{5}}{左は電話：1対1の配線。右は放送：1つのメッセージを周りの全員に一度に。}

ここまでニューロンは電話で話してきた。配線から配線へ、宛先を決めて。今度は最初の放送局をオンにしよう。セロトニンを放出する\nt{SERO}ニューロンである。脳全体で数十万個しかないが、その配線は脳の広大な領域に広がっている。

\section*{4つの違い}

\nt{SERO}ニューロンは\nt{GLUT}とは似ていない。第一に、その信号の符号は受け手が選ぶ。セロトニンには、どちらの符号の「錠」もある。第二に、自分で動く。覚醒中、一定の背景入力を受けていれば、毎秒数回の規則正しいカチッを鳴らす。メトロノームだ。第三に、遅い。その信号はミリ秒ではなく数百ミリ秒かけて効く。第四に、物質を細いすき間ではなく周囲の空間に放出することが多く、物質は水に落としたインクのしずくのように広がる。受け手は濃度を感じるが、分子が誰から来たのかはわからない。

\section*{存在しない論理}

抑制性の「錠」をもつメトロノームは、そのまま「NOT」になる。セロトニンがない間はニューロンが働き、現れると黙る。回路まるごとの代わりにニューロン1個。セロトニンのネットワークは、\nt{GLUT}ネットワークより論理的に強力に見える。

だが落とし穴がある。空間に放出された物質は混ざり合う。2つの信号が別々のままでいられるのは、物質が広がる距離、つまり数十マイクロメートルより離れた場所で放出されたときだけだ。ところがセロトニンニューロンの1個1個には、広い領域にばらまかれた膨大な数の放出部位があり、異なるニューロンの「配線」は重なり合う。こうしたネットワークで個々のビットは送れない。送れるのは少数の全体量だけだ。みんなが同じ局を聞いているラジオのようなものである。

\pencilfig{123}{%
% two release sites: far apart the clouds stay separate (two signals), close together they merge (one level)
\foreach \x/\y in {0.9/1.55, 4.1/1.55, 1.9/0, 2.9/0}{
  \foreach \r/\o in {0.95/0.07, 0.7/0.09, 0.45/0.12}{\fill[pcViolet, opacity=\o] (\x,\y) circle (\r);}
  \draw[dotpen=pcViolet, fill=pcViolet] (\x,\y) circle (0.07);}
\draw[arr] (1.95,1.55) -- (3.05,1.55); \draw[arr] (3.05,1.55) -- (1.95,1.55);
\node[lbl, anchor=south, align=center] at (2.5,1.6) {広がる距離\\より遠い};
\node[lbl, anchor=west] at (5.25,1.55) {2つの信号};
\node[lbl, anchor=west] at (5.25,0) {1つの共通の雲};
}{セロトニンは空間に放出され、雲のように広がる。2つの源が別々の信号を出せるのは、物質が広がる距離より互いに離れているときだけだ。近ければ雲は融合し、1つの共通レベルになる。}

\section*{実際に何をしているか}

1つの全体量を扱うための道具は、このシステムにそろっている。セロトニン濃度は個々のカチッをならして滑らかなレベルにする。最も単純な変換器が、スパイク列を電圧に変えるのと同じだ。このレベルは数秒続き、自然に消えていく。消去が必要なビットではなく、ゆっくり消える痕跡である。

さらにセロトニンニューロンには、自分自身に対する「錠」もある。放出されたセロトニンは自分の発生源にブレーキをかける。エンジニアなら、そこに負帰還増幅器、つまりレベルを保つサーモスタットを見てとるだろう。ただし実験によれば、この自己抑制は宛先を決めて働き、短い休止しか生まない。だからサーモスタットの役割は、今のところ事実というより仮説である。

\section*{待つことの値段}

ネットワーク全体で、ゆっくりと、しかも同じに保つ価値がある全体量とは何だろうか。有力な候補は\emph{忍耐}である。選択を想像してほしい。今の小さな報酬か、数秒後の大きな報酬か。どれだけ待てるかは、あなたにとって未来がどれだけ速く「値下がり」するかで決まる。ある仮説によれば、このつまみを回しているのがセロトニンだ。この仮説を支持する実験もある。セロトニンニューロンを人工的にオンにすると、マウスはあきらめるまでに、遅れて来る報酬をより長く待つ。

この章では、脳のすべての放送局が使う3つの装置を紹介した。メトロノーム型ニューロン、実は体積である「配線」、そして個々のカチッに代わるゆっくりしたレベルである。

\fullversion{5}
